\section{APPENDIX D}

\subsection{Paragraf 1: Justifikasi Bilah Galat}
\begin{frame}{Metodologi Estimasi}
    \begin{itemize}
        \item \textbf{Tujuan:} Justifikasi teoretis asal-usul \textit{error bars} pada visualisasi data.
        \item \textbf{Galat Populasi:} Estimasi berbasis galat eksperimental 8\% per gerbang 2-qubit.
        \item \textbf{Galat Statistik:} Kalkulasi galat 11\% berdasarkan jumlah tembakan (\textit{shots}) 8192 kali.
        \item \textbf{Total Galat:} Penulis mengasumsikan total galat 35\% untuk skenario terburuk.
        \item \textbf{Implementasi:} Penggunaan nilai unit counts ($\approx 360$) sebagai dasar bilah galat grafis.
    \end{itemize}
\end{frame}

\begin{frame}{Pertanyaan yang dijawab pada Paragraf 1}
    \centering
    \small
    \begin{tabular}{lp{4.5cm}p{6cm}}
        \toprule
        Kalimat ke & Pertanyaan & Jawaban \\
        \midrule
        1 & Tujuan lampiran? & Jelaskan asal-usul bilah galat grafis Gambar 2. \\
        4 & Galat statistik murni? & Sekitar 11\% untuk 8192 shots. \\
        5 & Estimasi total galat? & Diasumsikan 35\% (\textit{worst-case scenario}). \\
        7 & Nilai numerik per state? & Sekitar 360 unit counts. \\
        \bottomrule
    \end{tabular}
\end{frame}
